% ═══════════════════════════════════════════════════════════════════
%  Chapitre 4 — Sécurité de la plate-forme
% ═══════════════════════════════════════════════════════════════════
\chapter{Sécurité de la plate-forme}

\section*{Introduction}
\addcontentsline{toc}{section}{Introduction}

Ce chapitre est consacré à la sécurité de la solution. Il occupe une place
particulière dans ce rapport, et pour deux raisons.

La première tient à la nature du système. La plate-forme traite des données
personnelles --- identité, coordonnées, historique de fréquentation d'un
institut de beauté --- et porte des opérations à valeur économique :
réservations, commandes, points de fidélité convertibles en prestations
gratuites. Chacune de ces trois catégories appelle des garanties différentes,
et aucune n'est satisfaite par défaut.

La seconde tient au fait qu'un système exposé sur Internet n'a pas seulement
des utilisatrices : il a des visiteurs. L'API est publique par construction ---
l'application mobile l'appelle depuis les téléphones de clientes, sur des
réseaux quelconques. Tout ce qui n'est pas explicitement refusé côté serveur
est, en pratique, autorisé.

La démarche suivie n'a donc pas consisté à ajouter de la sécurité à un système
achevé, mais à traiter chaque décision de conception comme portant une part de
sécurité. Ce chapitre restitue cette démarche : le modèle de menaces retenu,
les principes directeurs, puis les contre-mesures effectivement implémentées,
domaine par domaine. Il se termine par une synthèse et par un inventaire
honnête des limites qui subsistent.

% ─────────────────────────────────────────────────────────────────
\section{Démarche et modèle de menaces}

\subsection{Surface d'attaque}

La première étape a consisté à délimiter la surface exposée et à situer les
frontières de confiance du système.

\figProjet[0.9]{securite-surface.png}{Surface d'attaque et frontières de confiance}{securite-surface}

Cinq points d'entrée existent, et un seul est réellement de confiance :

\begin{itemize}[leftmargin=1.2cm]
  \item \textbf{L'API HTTP}, atteignable par quiconque connaît son adresse. Sa
  documentation OpenAPI décrit d'ailleurs chacune de ses routes --- raison pour
  laquelle elle est coupée en production.

  \item \textbf{L'application mobile}, qui s'exécute sur le téléphone de la
  cliente. Elle échappe totalement à notre contrôle : son code est lisible, son
  trafic interceptable, ses contrôles contournables. \emph{Aucune règle de
  sécurité ne peut y résider.}

  \item \textbf{Le back-office web}, dans le même cas : c'est du JavaScript
  servi au navigateur. Masquer un bouton d'administration n'a jamais protégé la
  route correspondante.

  \item \textbf{Les images téléversées}, seuls fichiers que l'application écrit
  sur le disque du serveur --- et qu'elle republie ensuite en HTTP.

  \item \textbf{La base de données}, qui n'est pas exposée : son port n'est
  publié que sur l'interface locale du serveur.
\end{itemize}

Il en découle un principe qui structure tout ce chapitre :
\emph{la frontière de confiance passe à l'entrée de l'API, et nulle part
ailleurs}. Tout ce qui arrive d'un client --- corps de requête, en-têtes, type
MIME déclaré, extension de fichier, rôle affiché dans l'interface --- est une
donnée hostile jusqu'à vérification par le serveur.

\subsection{Modèle de menaces}

Les menaces ont été recensées selon les six catégories du modèle STRIDE, en
retenant pour chacune les scénarios réalistes pour ce système.

\begin{table}[H]
\centering
\caption{Modèle de menaces STRIDE appliqué à la plate-forme}
\label{tab:stride}
\begin{tabularx}{\textwidth}{L{2.9cm} X}
\toprule
\entete{Catégorie} & \entete{Scénarios retenus} \\
\midrule
\textbf{Usurpation}\newline\small(\textit{Spoofing}) &
Force brute sur le formulaire de connexion ; vol d'un jeton de session ;
devinette d'un code de réinitialisation à six chiffres ; forge d'un jeton
d'administration si le secret de signature est faible. \\

\textbf{Altération}\newline\small(\textit{Tampering}) &
Injection SQL ; modification du rôle porté dans la requête ; téléversement
d'un fichier exécutable déguisé en image ; falsification de l'en-tête
d'adresse IP réelle. \\

\textbf{Répudiation} &
Ajout manuel de points de fidélité sans trace ; remise d'un bon de récompense
sans enregistrement de qui l'a honoré. \\

\textbf{Divulgation}\newline\small(\textit{Information disclosure}) &
Lecture des commandes ou des rendez-vous d'une autre cliente ; fuite du schéma
de base par les messages d'erreur ; énumération des comptes existants ; export
massif de la base clientes. \\

\textbf{Déni de service} &
Saturation de l'API par un client unique ; épuisement du processeur par le
hachage de mots de passe ; export d'un volume de lignes qui bloque le
serveur ; réservation en masse de créneaux. \\

\textbf{Élévation de privilège} &
Accès direct à une route d'administration depuis un compte cliente ;
promotion de son propre rôle ; usage d'un jeton émis avant la suppression du
compte. \\
\bottomrule
\end{tabularx}
\end{table}

\subsection{Principes directeurs}

Quatre principes ont guidé les décisions.

\paragraph{Le serveur fait autorité.} Toute règle de sécurité est appliquée par
l'API. Les interfaces peuvent masquer, désactiver ou pré-valider : ce n'est
jamais qu'un confort d'usage, jamais une protection.

\paragraph{Défense en profondeur.} Aucune contre-mesure unique n'est réputée
suffisante. Un compte supprimé est ainsi protégé par trois verrous
indépendants : le refus de connexion sur son marqueur de suppression, le
remplacement de son mot de passe par le haché d'un secret que personne ne
connaît, et le rejet, à chaque requête, des jetons déjà émis.

\paragraph{Sécurité par défaut, et échec fermé.} Toute route est authentifiée
et comptée sauf déclaration contraire explicite. Toute propriété non déclarée
dans le contrat d'une route fait rejeter la requête, au lieu d'être ignorée en
silence. Et lorsqu'une configuration est incohérente, l'API \emph{refuse de
démarrer} plutôt que de fonctionner dans un état dégradé que rien ne
signalerait.

\paragraph{Moindre privilège.} Trois rôles seulement, chacun strictement borné.
L'image de production ne contient aucun compilateur et le processus n'y tourne
pas en tant que \texttt{root}. Le port de la base de données n'est pas publié.

% ─────────────────────────────────────────────────────────────────
\section{Authentification}

\subsection{Stockage des mots de passe}

Les mots de passe sont hachés avec \textbf{Argon2id}, lauréat du
\textit{Password Hashing Competition} et normalisé par la RFC 9106. Le choix
s'est porté sur lui plutôt que sur bcrypt pour sa résistance aux attaques
matérielles : Argon2 est conçu pour être coûteux en \emph{mémoire} autant qu'en
temps de calcul, ce qui ôte aux cartes graphiques et aux circuits dédiés
l'avantage massif dont ils disposent face aux fonctions purement calculatoires.

Aucun mot de passe n'est stocké, journalisé ou renvoyé en clair, à aucun
moment. La politique appliquée impose au minimum huit caractères, dont une
majuscule et une minuscule, et elle est \emph{la même} à l'inscription, au
changement de mot de passe et à la réinitialisation --- une exigence relâchée
sur l'un de ces trois chemins annulerait les deux autres.

\subsection{Jetons d'accès}

L'authentification des requêtes repose sur des jetons JWT signés, de durée
volontairement courte : \textbf{quinze minutes}. Cette brièveté est la réponse
au fait qu'un JWT est, par nature, invérifiable auprès d'une liste de
révocation sans requête supplémentaire : elle borne la fenêtre pendant laquelle
un jeton dérobé reste exploitable.

Le secret de signature est lu par une méthode qui \emph{lève une exception s'il
est absent}, aussi bien à la signature qu'à la vérification. C'est délibéré :
une valeur de repli en dur --- le réflexe habituel pour « que ça démarre en
développement » --- rendrait tous les jetons forgeables par quiconque a lu le
dépôt.

Le jeton n'est cependant pas cru sur parole. À chaque requête, la stratégie de
vérification recharge le compte en base et refuse le jeton si le compte
n'existe plus ou a été anonymisé.

\begin{lstlisting}[caption={Vérification du jeton d'accès : la signature ne suffit pas},captionpos=b]
async validate(payload: { sub: string; email: string; role: string }) {
  const user = await this.usersService.findById(payload.sub);
  // Un compte anonymise est traite comme inexistant. Sans ce test, l'access
  // token emis juste avant la suppression continuerait d'ouvrir le compte
  // jusqu'a son expiration.
  if (!user || user.deletedAt) {
    throw new UnauthorizedException();
  }
  return { id: user.id, email: user.email, role: user.role };
}
\end{lstlisting}

Ce rechargement a un second effet, moins évident et tout aussi important : le
rôle appliqué est celui \emph{de la base}, et non celui inscrit dans le jeton.
Une rétrogradation prend donc effet immédiatement, sans attendre l'expiration
des jetons en circulation.

\subsection{Jetons de rafraîchissement, rotation et détection de vol}

Des jetons d'accès de quinze minutes obligeraient la cliente à se reconnecter
plusieurs fois par jour. Un jeton de rafraîchissement de longue durée résout ce
problème --- et en crée un autre : c'est désormais lui qui vaut la session, et
sa durée de vie le rend d'autant plus intéressant à dérober.

Trois mesures répondent à ce risque.

\begin{itemize}[leftmargin=1.2cm]
  \item \textbf{Il n'est pas stocké en clair.} La base ne contient que son
  empreinte SHA-256. Une lecture de la base ne livre donc aucune session
  utilisable. Le hachage rapide suffit ici, contrairement au cas des mots de
  passe : le jeton est une valeur aléatoire de 48 octets, et non un secret
  choisi par un être humain --- il n'existe pas de dictionnaire de jetons.

  \item \textbf{Il tourne.} Chaque rafraîchissement révoque le jeton présenté
  et en émet un nouveau, rattaché à la même \emph{famille}.

  \item \textbf{Sa réutilisation est détectée.} Présenter deux fois le même
  jeton est impossible en usage normal. C'est en revanche exactement ce qui se
  produit lorsqu'un jeton a été dérobé : l'attaquant et la cliente légitime
  l'utilisent chacun de leur côté. Impossible de savoir lequel des deux se
  présente --- le système révoque donc \emph{toute la famille}. Les deux sont
  déconnectés, et l'accès frauduleux prend fin. Le coût pour la cliente est une
  reconnexion ; le bénéfice est la fin d'une session volée qui, sans cela, se
  prolongerait indéfiniment par rotations successives.
\end{itemize}

Côté mobile, le jeton n'est pas rangé dans le stockage ordinaire de
l'application mais dans le \textbf{coffre sécurisé du système
d'exploitation} --- trousseau iOS, \textit{Keystore} Android --- adossé au
matériel du téléphone.

Le déroulement complet, rotation et détection de réutilisation comprises, est
donné par la \figref{seq-refresh} du chapitre~2.

\subsection{Codes à usage unique}

La vérification d'adresse et la réinitialisation de mot de passe reposent sur
un code à six chiffres envoyé par courriel. Un tel code est court par
nécessité --- il doit être recopié à la main --- donc faible par construction :
un million de possibilités se parcourent en quelques minutes. Cinq mesures le
rendent néanmoins sûr.

\begin{table}[H]
\centering
\caption{Protection des codes à usage unique}
\label{tab:codes}
\begin{tabularx}{\textwidth}{L{4.0cm} X}
\toprule
\entete{Mesure} & \entete{Ce qu'elle empêche} \\
\midrule
Génération par \texttt{randomInt} &
Générateur cryptographique, et non \texttt{Math.random} dont la suite est
prédictible : un code observé ne permet pas de deviner le suivant. \\

Stockage haché &
Une lecture de la base ne livre aucun code exploitable. \\

Plafond de cinq tentatives &
Ferme la porte à l'énumération : sans lui, six chiffres se devinent. \\

Validité de quinze minutes &
Borne la fenêtre d'exploitation d'un code intercepté. \\

Comparaison en temps constant &
\texttt{timingSafeEqual} : le temps de réponse ne trahit pas le nombre de
chiffres corrects, ce qui interdit de reconstituer le code position par
position. \\
\bottomrule
\end{tabularx}
\end{table}

À quoi s'ajoute une propriété structurelle : la recherche s'effectue par
\emph{couple (compte, usage)} et jamais par le code seul. Un code deviné ne
vaut donc que pour le compte auquel il a été émis --- alors qu'une recherche
par code seul aurait fait qu'un code trouvé au hasard ouvre n'importe quel
compte l'ayant reçu, ce qui multiplie les chances de l'attaquant par le nombre
de codes en circulation. Enfin, émettre un nouveau code efface le précédent du
même usage : au plus un code actif par cliente et par usage.

\subsection{Défense contre l'énumération de comptes}

Savoir quelles adresses possèdent un compte est une information exploitable :
elle alimente le harponnage ciblé et, s'agissant d'un institut de beauté, elle
révèle une information personnelle en soi.

Deux canaux la trahissent, et tous deux ont été fermés.

\paragraph{Le canal des messages.} La demande de réinitialisation répond la
même chose que l'adresse existe ou non.

\paragraph{Le canal du temps.} Ce canal est plus subtil, et c'est précisément
ce qui le rend dangereux : il subsiste alors même que les messages sont
identiques. Vérifier un mot de passe avec Argon2 coûte plusieurs centaines de
millisecondes. Si le serveur répondait immédiatement pour une adresse inconnue
--- puisqu'il n'a alors aucun haché à comparer --- et après ce calcul pour une
adresse connue, la seule mesure du délai de réponse suffirait à trier les
adresses.

La connexion exécute donc, lorsque l'adresse est inconnue, une vérification
\emph{factice} contre le haché d'une valeur aléatoire. Elle échoue toujours :
seul le temps consommé compte.

\begin{lstlisting}[caption={Uniformisation du temps de réponse à la connexion},captionpos=b]
// Hash de reference, calcule une seule fois, sur une valeur aleatoire que
// personne ne connait. Il ne protege rien : il sert uniquement a faire
// travailler argon2 aussi longtemps qu'une vraie verification, pour qu'un
// email inconnu reponde au meme rythme.
private async burnPasswordComparison(candidate: string): Promise<void> {
  this.dummyHashPromise ??= argon2.hash(randomBytes(32).toString('hex'));
  const dummyHash = await this.dummyHashPromise;
  await argon2.verify(dummyHash, candidate).catch(() => false);
}
\end{lstlisting}

% ─────────────────────────────────────────────────────────────────
\section{Autorisation et contrôle d'accès}

\subsection{Modèle de contrôle d'accès}

Le contrôle d'accès repose sur un modèle par rôles (RBAC) à trois niveaux
strictement ordonnés. Le tableau ci-dessous donne la matrice des droits
effectivement appliquée par le serveur.

\begin{table}[H]
\centering
\caption{Matrice des droits par rôle}
\label{tab:matrice-droits}
\begin{tabularx}{\textwidth}{X C{1.6cm} C{1.7cm} C{2.1cm} C{1.8cm}}
\toprule
\entete{Opération} & \entete{Public} & \entete{Cliente} & \entete{Personnel} & \entete{Gérante} \\
\midrule
Consulter les catalogues actifs   & \checkmark & \checkmark & \checkmark & \checkmark \\
Réserver, commander pour soi      &            & \checkmark & \checkmark & \checkmark \\
Consulter \emph{ses} rendez-vous et commandes & & \checkmark & \checkmark & \checkmark \\
Consulter ceux de \emph{toutes}   &            &            & \checkmark & \checkmark \\
Réserver pour le compte d'une cliente &        &            & \checkmark & \checkmark \\
Changer un statut, honorer un bon &            &            & \checkmark & \checkmark \\
Créditer des points à la main     &            &            & \checkmark & \checkmark \\
Gérer le catalogue et les stocks  &            &            &            & \checkmark \\
Régler horaires, capacité, fermetures &        &            &            & \checkmark \\
Définir récompenses et paliers    &            &            &            & \checkmark \\
Gérer les comptes et les rôles    &            &            &            & \checkmark \\
Consulter le journal d'audit      &            &            &            & \checkmark \\
Exporter les données (Excel)      &            &            &            & \checkmark \\
Téléverser une image              &            &            &            & \checkmark \\
\bottomrule
\end{tabularx}
\end{table}

\subsection{Application par gardes successifs}

Cette matrice n'est pas une convention documentaire : elle est appliquée par
deux gardes que toute requête traverse, au sein d'une chaîne de filtres qu'elle
parcourt \emph{avant} d'atteindre la moindre ligne de logique métier.

\figProjet[0.85]{securite.png}{Chaîne de traitement de sécurité d'une requête}{securite}

Le premier vérifie la signature et la validité du jeton. Il est actif par
défaut : une route n'échappe à l'authentification que si elle est
\emph{explicitement} déclarée publique. Ce sens du défaut est essentiel ---
avec la convention inverse, une route ajoutée un jour de fatigue serait
ouverte, et rien ne le signalerait.

Le second compare le rôle du compte aux rôles exigés par la route. Il est
indépendant de tout ce que l'interface affiche ou masque : une route
d'administration atteinte par un autre moyen --- l'outil de développement du
navigateur, un client HTTP en ligne de commande, la documentation OpenAPI ---
est refusée par le serveur lui-même.

\subsection{Contrôle de propriété des ressources}

Le rôle ne suffit pas. Deux clientes portent le même rôle, et rien dans ce rôle
n'empêche l'une de demander la commande de l'autre : c'est la faille de type
\textit{Insecure Direct Object Reference}, l'une des plus répandues et des plus
faciles à exploiter --- il suffit de changer un identifiant dans une adresse.

Chaque route qui expose une ressource nominative vérifie donc, en plus du rôle,
que la ressource appartient bien au demandeur. La vérification est faite
\emph{à la lecture en base}, en filtrant sur le propriétaire, et non après
coup : une commande qui n'appartient pas à la cliente est traitée comme
introuvable, ce qui évite au passage de confirmer son existence.

Les identifiants sont par ailleurs des UUID et non des entiers séquentiels. Ce
n'est pas une protection en soi --- la vérification de propriété reste la seule
qui compte --- mais cela ôte à un attaquant la possibilité d'énumérer les
ressources en incrémentant un nombre, et empêche de déduire le volume
d'activité de l'institut d'un identifiant observé.

\subsection{Invariants d'administration}

Deux invariants protègent le compte de direction, et ils sont appliqués par le
serveur :

\begin{itemize}[leftmargin=1.2cm]
  \item \textbf{une seule administratrice à la fois} --- promouvoir un compte
  exige de vérifier qu'aucun autre ne porte déjà ce rôle ;
  \item \textbf{jamais zéro administratrice} --- rétrograder la dernière est
  refusé. Sans ce garde-fou, une fausse manipulation rendrait la gestion du
  centre définitivement inaccessible, sans autre recours qu'une intervention
  directe en base.
\end{itemize}

% ─────────────────────────────────────────────────────────────────
\section{Sécurité des entrées et des sorties}

\subsection{Validation stricte des entrées}

Un tube de validation global s'applique à toutes les routes, avec trois
réglages qui comptent :

\begin{itemize}[leftmargin=1.2cm]
  \item \textbf{liste blanche} --- seules les propriétés déclarées dans le
  contrat de la route sont conservées ;
  \item \textbf{rejet des propriétés inconnues} --- une propriété non déclarée
  fait \emph{échouer} la requête au lieu d'être silencieusement ignorée. C'est
  ce réglage qui ferme la porte à l'affectation en masse : soumettre un champ
  \texttt{role} ou \texttt{pointsBalance} dans une mise à jour de profil est
  rejeté, et non toléré ;
  \item \textbf{conversion typée} --- les valeurs sont converties dans le type
  attendu avant d'atteindre la logique métier, qui ne reçoit donc jamais une
  chaîne là où elle attend un nombre ou une date.
\end{itemize}

Les règles de validation sont déclarées au plus près des données, dans les
objets de transfert, et sont partagées : la politique de mot de passe, le
format des codes à six chiffres et celui des numéros de téléphone marocains
sont des validateurs uniques, réutilisés partout où la donnée apparaît.

\subsection{Injection SQL}

Toutes les requêtes passent par l'ORM, qui les émet sous forme paramétrée :
les valeurs ne sont jamais concaténées dans le texte de la requête.

Le projet compte \emph{une} requête SQL écrite à la main --- la prise du verrou
de réservation. Elle est écrite avec un gabarit étiqueté (\textit{tagged
template}), lui aussi paramétré, et sa seule valeur variable est une constante
du code, jamais une entrée utilisateur.

\subsection{Contrôle des fichiers téléversés}

Le téléversement d'images est le seul chemin par lequel un utilisateur écrit un
fichier sur le serveur --- fichier ensuite republié en HTTP. C'est donc le
point le plus sensible du système, et il est traité comme tel.

\begin{itemize}[leftmargin=1.2cm]
  \item \textbf{Réservé à la gérante.} Aucune cliente ne peut téléverser quoi
  que ce soit.

  \item \textbf{Reconnaissance par signature binaire.} Le type déclaré par le
  navigateur et l'extension du nom de fichier viennent tous deux du client et
  se falsifient en une ligne. Seuls les premiers octets du fichier disent ce
  qu'il \emph{est} : trois signatures sont acceptées --- JPEG, PNG, WebP ---
  et c'est la signature reconnue, jamais l'extension fournie, qui détermine
  l'extension enregistrée.

  \item \textbf{Le SVG est refusé.} C'est un refus délibéré, et non un oubli :
  le SVG est du XML susceptible de contenir du JavaScript. Servi depuis le
  domaine de l'application, il s'exécuterait avec les droits de celle-ci ---
  une faille XSS stockée, ouverte par un simple téléversement d'image.

  \item \textbf{Nom de fichier tiré au sort.} Le nom fourni par le client est
  jeté ; le fichier est enregistré sous un identifiant aléatoire. Cela ferme la
  traversée de répertoire (un nom du type \texttt{../../etc/passwd}) et
  l'écrasement d'un fichier existant.

  \item \textbf{Taille plafonnée à 5 Mo}, et un seul fichier par requête.

  \item \textbf{Cinq minutes de recul.} Les images sont servies avec un cache
  d'un an, marqué immuable --- ce qui est sûr précisément parce que le nom est
  tiré au sort : remplacer une photo crée un nouveau nom, jamais une nouvelle
  version du même. Il n'y a donc rien à invalider.
\end{itemize}

\subsection{Fuite d'information par les messages d'erreur}

Un message d'erreur technique renseigne un attaquant : noms de tables, noms de
colonnes, noms de contraintes, parfois valeurs. Un filtre global intercepte
donc toute exception avant qu'elle n'atteigne le client et applique trois
traitements :

\begin{enumerate}[leftmargin=1.2cm]
  \item les exceptions métier levées volontairement passent intactes --- leur
  message est destiné à l'utilisatrice ;
  \item les erreurs de base de données identifiées sont traduites en langage
  métier ;
  \item \textbf{tout le reste devient un message neutre}, le détail et la pile
  d'appels ne vivant que dans les journaux du serveur.
\end{enumerate}

Les journaux eux-mêmes ont été traités comme un risque : une erreur serveur y
consigne l'\emph{identifiant} de l'utilisateur, jamais son adresse
électronique --- assez pour rejouer l'incident, pas assez pour reconstituer un
fichier d'adresses à partir des journaux.

\subsection{Injection d'en-tête de courriel}

Le nom de l'expéditeur est inséré entre guillemets dans un en-tête
\texttt{From}. Un retour à la ligne y couperait l'en-tête en deux et
permettrait d'en injecter d'autres --- un \texttt{Bcc} vers une adresse
tierce, par exemple. Les caractères de rupture sont donc retirés avant
insertion. La valeur vient certes de la configuration et non d'un client, mais
un fichier d'environnement recopié de travers suffirait à produire l'effet, et
l'échec serait silencieux.

% ─────────────────────────────────────────────────────────────────
\section{Sécurité du transport et des en-têtes}

\subsection{Chiffrement du transport}

L'API n'est publiée que sur l'interface locale du serveur. Ce choix est une
mesure de sécurité et non une limitation : il rend \emph{impossible} de la
joindre sans passer par le reverse proxy placé devant, à qui revient la
terminaison HTTPS. Publier le port directement reviendrait à servir en clair
les jetons d'authentification, sur des réseaux Wi-Fi publics compris.

Côté mobile, un garde-fou interdit de configurer une adresse d'API non
chiffrée en dehors du développement local.

\subsection{En-têtes de sécurité HTTP}

Helmet retire l'en-tête qui annonçait la technologie du serveur --- une
information sans usage légitime, qui oriente le choix des exploits --- et pose
une douzaine d'en-têtes que le navigateur applique.

\begin{table}[H]
\centering
\caption{Principaux en-têtes de sécurité posés}
\label{tab:entetes}
\begin{tabularx}{\textwidth}{L{6.3cm} X}
\toprule
\entete{En-tête} & \entete{Rôle} \\
\midrule
\texttt{Strict-Transport-Security} & Impose HTTPS pour les visites suivantes,
et ferme la fenêtre de rétrogradation vers HTTP. \\
\texttt{Content-Security-Policy} & Restreint les sources de scripts : première
barrière contre le XSS. \\
\texttt{X-Content-Type-Options} & Interdit au navigateur de deviner le type
d'une ressource, donc d'exécuter comme script un fichier qui n'en est pas un. \\
\texttt{X-Frame-Options} & Interdit l'inclusion dans un cadre : protection
contre le détournement de clic. \\
\texttt{Cross-Origin-Resource-Policy} & Contrôle quelles origines peuvent
charger une ressource. \\
\bottomrule
\end{tabularx}
\end{table}

\begin{attention}[Un arbitrage assumé]
La politique de sécurité du contenu interdit par défaut les scripts en ligne.
Or l'interface de documentation OpenAPI s'initialise justement par un script en
ligne : l'activer telle quelle afficherait une page blanche. Elle est donc
désactivée \emph{uniquement là où cette documentation tourne}, c'est-à-dire en
développement. En production, la documentation est coupée et la politique
s'applique pleinement. L'arbitrage porte ainsi sur un environnement qui n'est
pas exposé, et non sur le serveur en ligne.
\end{attention}

\subsection{Contrôle des origines}

La politique CORS n'autorise que les adresses déclarées du back-office. Un site
tiers ne peut donc pas faire émettre à un navigateur des requêtes authentifiées
vers l'API.

\begin{attention}[Un piège de déploiement]
Cette liste doit contenir l'adresse réelle du back-office déployé. Non
renseignée, elle retombe sur l'adresse de développement : le back-office ne
peut plus joindre l'API, et l'échec ne se manifeste que par un blocage dans la
console du navigateur --- \emph{jamais} par un message du serveur. On cherche
donc longtemps du mauvais côté.
\end{attention}

% ─────────────────────────────────────────────────────────────────
\section{Disponibilité et résistance à l'abus}

\subsection{Limitation du débit}

Un compteur par adresse IP plafonne l'ensemble de l'API à 120 requêtes par
minute. Les routes coûteuses --- celles qui déclenchent un hachage Argon2 ou un
envoi de courriel --- portent en plus une limite stricte.

\begin{table}[H]
\centering
\caption{Limites de débit par catégorie de route}
\label{tab:rate-limit}
\begin{tabularx}{\textwidth}{L{5.4cm} C{3.0cm} X}
\toprule
\entete{Route} & \entete{Limite} & \entete{Ce qu'elle protège} \\
\midrule
Ensemble de l'API & 120 / min &
Saturation générale du service \\
Connexion, inscription & 10 / min &
Force brute sur les mots de passe ; épuisement du processeur par Argon2 \\
Changement de mot de passe & 10 / min & Force brute sur le mot de passe actuel \\
Vérification d'adresse & 10 / min & Énumération du code à six chiffres \\
Renvoi de code & 5 / min & Usage de l'API comme relais d'envoi de courriels \\
Mot de passe oublié & 5 / 5 min & Harcèlement par courriel d'une adresse tierce \\
Téléversement d'image & 30 / min & Saturation du disque du serveur \\
\bottomrule
\end{tabularx}
\end{table}

\begin{attention}[Le réglage qui décide de tout]
Le compteur s'appuie sur l'adresse IP du demandeur. Derrière un reverse proxy,
cette adresse est celle \emph{du proxy}, identique pour tout le monde : sans le
réglage correspondant, l'institut entier partage un unique compteur de 120
requêtes par minute, et l'API se met à répondre 429 au hasard dès qu'il y a du
monde. Le symptôme est déroutant --- parfait en test avec une seule personne,
cassé en production --- et rien dans les journaux ne le rattache à sa cause.

Le réglage inverse est tout aussi dangereux, et c'est ce qui en fait un vrai
arbitrage : activé sans proxy devant, l'en-tête d'adresse réelle devient
falsifiable par n'importe qui, et il suffit alors d'en changer la valeur à
chaque requête pour contourner entièrement la limitation. Ce réglage n'est donc
à activer \emph{que} si un proxy de confiance est réellement en place.
\end{attention}

\subsection{Plafonnement des volumes}

Deux mesures évitent qu'une requête légitime ne dégénère en déni de service.

\paragraph{La pagination obligatoire.} Toute liste susceptible de croître ---
historique de fidélité, commandes, rendez-vous, utilisateurs --- est paginée
côté serveur. Sans cela, une cliente fidèle depuis trois ans rechargeait la
totalité de son historique à chaque ouverture de l'écran.

\paragraph{Le plafond d'export.} Les exports Excel comptent les lignes
\emph{avant} de charger quoi que ce soit et refusent au-delà de 20 000 lignes,
en invitant à restreindre la période. Un export non borné aurait constitué
l'opération la plus lourde de l'API, accessible d'une seule requête. Le
comptage préalable importe autant que le plafond : c'est lui qui garantit que
le refus ne coûte rien.

\subsection{Correction sous concurrence}

Les mécanismes décrits au chapitre~2 --- verrou consultatif sur les
réservations, écritures conditionnelles atomiques sur le stock et sur le solde
de points --- relèvent aussi de la sécurité, et pas seulement de la
correction fonctionnelle.

Une condition de course est en effet exploitable dès lors qu'elle est
déclenchable à volonté. Le débit conditionnel du solde de points en est
l'illustration : lire le solde puis le décrémenter permettrait, en lançant
plusieurs échanges simultanés, de dépenser plusieurs fois les mêmes points et
d'obtenir des récompenses non acquises. Le débit est donc exprimé comme une
\emph{écriture conditionnelle} --- retirer $n$ points à condition qu'il y en ait
au moins $n$ --- évaluée par la base au moment de l'écriture. Si la condition
n'est pas satisfaite, aucune ligne n'est modifiée et rien n'est émis.

% ─────────────────────────────────────────────────────────────────
\section{Traçabilité et non-répudiation}

Certaines opérations ont une valeur économique et doivent laisser une trace
opposable.

\begin{itemize}[leftmargin=1.2cm]
  \item \textbf{L'ajout manuel de points} exige un motif et enregistre
  l'employée qui l'a saisi. Un journal d'audit, consultable par la seule
  gérante, liste ces écritures.

  \item \textbf{La remise d'un bon} enregistre la date et l'identité de la
  personne qui l'a honoré.

  \item \textbf{Le lien vers l'employée n'est jamais supprimé en cascade.} Le
  départ d'une employée ne doit ni effacer les points des clientes qu'elle a
  servies, ni la trace de ce qui leur a été remis : l'écriture demeure, elle ne
  désigne simplement plus personne.

  \item \textbf{Les montants sont figés} au moment de l'acte. Une révision de
  tarif ne réécrit pas rétroactivement ce qui a été facturé, ni les points
  gagnés.
\end{itemize}

% ─────────────────────────────────────────────────────────────────
\section{Protection des données personnelles}

\subsection{Cadre applicable}

Le traitement de données personnelles au Maroc est régi par la \textbf{loi
09-08}, dont la Commission Nationale de contrôle de la protection des Données à
caractère Personnel (CNDP) assure le contrôle. Trois obligations ont eu un
effet direct sur la conception.

\paragraph{La minimisation.} Seules sont collectées les données nécessaires :
identité, adresse électronique, téléphone --- ce dernier parce qu'un institut
doit pouvoir rappeler une cliente pour confirmer une commande. Aucune donnée de
paiement n'est traitée : le règlement s'effectue à la remise, en espèces, et
l'application n'en voit rien. C'est une réduction considérable du risque, et
elle découle d'un choix fonctionnel plutôt que technique.

\paragraph{Le droit d'accès.} La cliente consulte à tout moment son profil, ses
rendez-vous, ses commandes et l'intégralité de son historique de fidélité.

\paragraph{Le droit à l'effacement.} C'est celui qui a demandé le plus de
travail, et il fait l'objet de la section suivante.

\subsection{L'effacement par anonymisation}

Le droit à l'effacement se heurte ici à une exigence contradictoire : les
commandes et les rendez-vous d'une cliente portent le chiffre d'affaires de
l'institut, qui doit lui survivre. Une suppression en cascade détruirait la
comptabilité ; un refus de supprimer méconnaîtrait la loi.

La solution retenue est l'\textbf{anonymisation} : la ligne survit, vidée de
tout ce qui identifie la personne.

\begin{lstlisting}[caption={Anonymisation d'un compte : ce qui est effacé, ce qui survit},captionpos=b]
await this.prisma.$transaction([
  this.prisma.user.update({
    where: { id: userId },
    data: {
      // L'adresse d'origine est LIBEREE : si la cliente revient un jour, elle
      // pourra se reinscrire avec le meme email — sur un compte neuf. Le
      // domaine `.invalid` est reserve par la RFC 2606 et n'est routable nulle
      // part : aucun email ne partira jamais vers cette coquille.
      email: `supprime-${userId}@compte-supprime.invalid`,
      firstName: null, lastName: null, phone: null,
      passwordHash,          // hache d'un secret aleatoire que nul ne connait
      emailVerifiedAt: null,
      deletedAt: new Date(),
    },
  }),
  // Les sessions ouvertes tombent tout de suite.
  this.prisma.refreshToken.deleteMany({ where: { userId } }),
  // Et un code de reinitialisation en cours ne doit pas servir de porte
  // derobee sur une coquille anonyme.
  this.prisma.verificationCode.deleteMany({ where: { userId } }),
]);
\end{lstlisting}

Plusieurs points de ce traitement méritent d'être relevés.

\begin{itemize}[leftmargin=1.2cm]
  \item \textbf{L'opération est transactionnelle.} L'anonymisation, la
  destruction des sessions et celle des codes en cours réussissent ou échouent
  ensemble. Un échec partiel laisserait une coquille anonyme accessible par une
  session encore ouverte.

  \item \textbf{Le domaine \texttt{.invalid} est réservé par la RFC 2606} et
  n'est routable nulle part. Aucun courriel ne peut partir vers cette coquille,
  même par erreur de code --- ce qu'un domaine inventé ne garantirait pas.

  \item \textbf{L'adresse d'origine est libérée}, ce qui permet une
  réinscription ultérieure sur un compte neuf.

  \item \textbf{Trois verrous indépendants} ferment le compte : le refus de
  connexion sur le marqueur de suppression, le mot de passe devenu inconnu de
  tous, et le rejet à chaque requête des jetons d'accès déjà émis.

  \item \textbf{La garantie est structurelle.} Les rendez-vous et les commandes
  sont rattachés à la cliente en mode \texttt{Restrict} : si du code tentait un
  jour une suppression, la contrainte le ferait échouer bruyamment au lieu de
  détruire l'historique en silence.
\end{itemize}

L'opération est enfin \textbf{idempotente} : rejouée sur un compte déjà
anonymisé, elle répond succès sans rien réécrire --- ce qui préserve la date
d'origine, seule trace de la date à laquelle le droit a été exercé.

\subsection{Gestion des secrets}

Aucun secret ne figure dans le dépôt. Les fichiers d'environnement ne sont pas
versionnés ; seuls leurs modèles le sont, documentés valeur par valeur.

La configuration est \textbf{validée au démarrage}, et l'API refuse de démarrer
si elle est incohérente. Ce choix mérite d'être souligné : la solution
alternative --- démarrer et se débrouiller --- produit des états dégradés que
rien ne signale.

\begin{itemize}[leftmargin=1.2cm]
  \item \textbf{Le secret de signature doit compter au moins 32 caractères.} Un
  secret plus court se casse hors ligne, et permet ensuite de forger un jeton
  d'administration sans jamais toucher à l'API.

  \item \textbf{L'envoi réel de courriels est exigé en production.} Un serveur
  de production laissé en mode console n'enverrait aucun code de vérification
  ni de réinitialisation, et rien ne le signalerait : les clientes se
  retrouveraient simplement incapables de valider leur compte.

  \item \textbf{Les valeurs de repli sont bannies} là où elles seraient
  dangereuses. Le secret de signature est lu par une méthode qui lève une
  exception s'il manque, à la signature comme à la vérification.
\end{itemize}

\subsection{Sauvegardes}

Un script produit un cliché horodaté de la base et supprime ceux de plus de
quatorze jours. Deux règles d'exploitation l'accompagnent, et elles comptent
autant que le script : le cliché doit être recopié \emph{hors de la machine}
--- une sauvegarde qui vit sur le serveur qu'elle sauvegarde ne protège de rien
--- et le dossier des images, qui vit hors de la base, doit être copié
séparément. La procédure de restauration et sa vérification sont documentées
dans le script lui-même.

% ─────────────────────────────────────────────────────────────────
\section{Sécurité de la chaîne de développement}

La sécurité d'une application ne s'arrête pas à son code : elle englobe ses
dépendances, sa construction et son déploiement.

\paragraph{Les dépendances.} Les versions sont figées par un fichier de
verrouillage, et deux dépendances \emph{transitives} sont forcées à des
versions corrigées : une bibliothèque d'analyse YAML tirée par la couche de
documentation, et une bibliothèque d'identifiants tirée par le générateur de
classeurs Excel. Ces dépendances-là sont les plus faciles à négliger --- on ne
les a pas choisies, et elles n'apparaissent pas dans la liste des dépendances
directes.

\paragraph{L'image de production.} Elle est construite en trois étapes, si bien
que l'image livrée ne contient ni compilateur, ni cadriciel de test, ni
outillage de développement --- une réduction de surface d'attaque autant qu'un
gain de taille. Le processus y tourne sous un utilisateur non privilégié, et
non en tant que \texttt{root}.

\paragraph{L'intégration continue.} Chaque modification rejoue l'intégralité des
tests contre une véritable base PostgreSQL. Plusieurs de ces tests portent
directement sur des propriétés de sécurité : rotation et détection de
réutilisation des jetons, anonymisation des comptes, capacité des créneaux sous
concurrence, restriction des exports.

\paragraph{L'hygiène du dépôt.} Les notes de travail internes, qui décrivaient
des faiblesses avant leur correction, ont été retirées du suivi. Le retrait est
documenté avec sa limite --- il ne réécrit pas l'historique, et le contenu
demeure accessible dans les commits passés. Cette honnêteté a une valeur
opérationnelle : elle évite de tenir pour effacé ce qui ne l'est pas.

% ─────────────────────────────────────────────────────────────────
\section{Synthèse}

Le tableau ci-dessous récapitule, pour chaque menace du modèle, la
contre-mesure retenue et l'endroit où elle est appliquée.

\begin{table}[H]
\centering
\caption{Synthèse menaces $\rightarrow$ contre-mesures}
\label{tab:synthese-securite}
\begin{tabularx}{\textwidth}{L{4.2cm} X}
\toprule
\entete{Menace} & \entete{Contre-mesure} \\
\midrule
Force brute sur les mots de passe &
Argon2id ; 10 tentatives par minute et par IP ; politique de complexité
identique sur les trois chemins de saisie \\

Vol d'un jeton de session &
Jeton d'accès de 15 minutes ; rafraîchissement haché, tourné, avec révocation
de toute la famille en cas de réutilisation ; coffre sécurisé du téléphone \\

Devinette d'un code à 6 chiffres &
Générateur cryptographique ; stockage haché ; 5 tentatives ; 15 minutes de
validité ; comparaison en temps constant ; recherche par (compte, usage) \\

Énumération de comptes &
Réponses identiques et \emph{temps de réponse} identiques, par vérification
factice sur adresse inconnue \\

Élévation de privilège &
Garde de rôle côté serveur sur chaque route ; rôle relu en base à chaque
requête ; rejet des propriétés non déclarées \\

Lecture des données d'autrui &
Contrôle de propriété au niveau de la requête en base ; identifiants UUID \\

Jeton émis avant suppression &
Compte rechargé et marqueur de suppression vérifié à chaque requête \\

Injection SQL &
Requêtes paramétrées par l'ORM ; unique requête manuelle sans entrée
utilisateur \\

XSS par téléversement &
Reconnaissance par signature binaire ; SVG refusé ; nom aléatoire ; en-têtes
\texttt{nosniff} et CSP \\

Fuite du schéma de base &
Filtre d'exceptions global ; documentation OpenAPI coupée en production \\

Déni de service &
Limitation de débit globale et par route ; pagination obligatoire ; plafond
d'export à 20 000 lignes \\

Condition de course exploitable &
Verrou consultatif sur les réservations ; écritures conditionnelles atomiques
sur le stock et le solde \\

Interception du trafic &
API sur interface locale ; HTTPS imposé par le reverse proxy ; HSTS ;
garde-fou HTTPS côté mobile \\

Répudiation d'une opération &
Motif obligatoire et auteur enregistré ; journal d'audit ; liens vers
l'employée jamais supprimés en cascade \\

Atteinte aux données personnelles &
Minimisation ; aucune donnée de paiement ; anonymisation transactionnelle ;
sauvegardes horodatées \\
\bottomrule
\end{tabularx}
\end{table}

% ─────────────────────────────────────────────────────────────────
\section{Limites connues et durcissement à venir}

Aucun système n'est sûr dans l'absolu, et un rapport qui prétendrait le
contraire manquerait son objet. Les limites suivantes sont assumées et
identifiées.

\paragraph{L'absence de second facteur.} L'authentification repose sur un seul
facteur. Pour un compte cliente, l'enjeu reste mesuré ; pour le compte de
direction, qui donne accès à l'ensemble du fichier clientes, un second facteur
serait la première amélioration à apporter.

\paragraph{Les données au repos ne sont pas chiffrées.} La base repose sur le
chiffrement du disque du serveur hôte. Un chiffrement au niveau des colonnes
les plus sensibles constituerait une couche supplémentaire.

\paragraph{Le compteur de débit vit en mémoire.} Il est donc propre à chaque
instance et remis à zéro au redémarrage. Tant que l'API tourne en un seul
exemplaire, la protection est effective ; une mise à l'échelle horizontale
imposerait de déporter ce compteur dans un magasin partagé.

\paragraph{Aucun audit externe n'a été mené.} Les propriétés de sécurité sont
établies par les tests automatisés et par la revue du code. Un test
d'intrusion mené par un tiers reste la seule façon de découvrir ce que la
personne qui a écrit le code ne pense pas à chercher.

\paragraph{La supervision est minimale.} Les événements sont journalisés, mais
rien ne les agrège ni ne déclenche d'alerte. Une série d'échecs de connexion
sur un même compte devrait alerter, et ne le fait pas aujourd'hui.

\paragraph{Les points restant à faire avant l'ouverture publique.} Le
certificat HTTPS et le reverse proxy conditionnent plusieurs des mesures
décrites ici : tant qu'ils ne sont pas en place, la protection du transport
repose sur le fait que le service n'est pas encore exposé.

\section*{Conclusion}
\addcontentsline{toc}{section}{Conclusion}

Ce chapitre a présenté la sécurité de la plate-forme comme une chaîne
continue, du modèle de menaces jusqu'aux limites qui subsistent. Trois idées
en résument la démarche.

D'abord, \textbf{la frontière de confiance passe à l'entrée de l'API}. Ni
l'application mobile, ni le back-office, ni aucune valeur venue d'un client
n'est cru sur parole ; ce qui est masqué dans une interface n'est jamais
protégé pour autant.

Ensuite, \textbf{les défauts sont fermés}. Une route est authentifiée et
comptée sauf déclaration contraire, une propriété inconnue fait échouer la
requête, et une configuration incohérente empêche le démarrage plutôt que de
produire un état dégradé silencieux.

Enfin, \textbf{la sécurité n'a pas été ajoutée après coup}. Les mécanismes qui
garantissent qu'un créneau n'est pas vendu deux fois sont les mêmes qui
empêchent de dépenser deux fois les mêmes points ; le figement des montants
sert à la fois l'exactitude comptable et la non-répudiation ; l'anonymisation
concilie un droit légal et la survie de l'historique. Correction fonctionnelle
et sécurité se sont révélées être, le plus souvent, deux vues du même
problème.

Le chapitre suivant présente la mise en \oe uvre effective de la solution et
les tests qui en établissent les propriétés.
